\chapter{Code and Reproducibility}
\label{app:code}

The complete source code, including the notebook, the raw and processed
data (tracked with Git LFS), and this report's own \LaTeX{} source, is
available at
\url{https://github.com/sepanta007/Matrix_Factorization}. The listings
below reproduce the core modules directly for convenience, exactly as used
to produce the results in this report.

\section{Repository layout}

The project repository is organized as follows; every path below is
relative to the project root.

\begin{lstlisting}[style=plaintext, caption={Repository layout.}]
src/
  eals.py             eALS core algorithm (Algorithm 1), PySpark RDD-based
  baselines.py        Spark MLlib ALS reference baseline
  evaluate.py         HR@K / NDCG@K leave-one-out evaluation
  preprocess.py        raw Amazon reviews -> filtered train/test parquet
  run_experiments.py    weighting / convergence / scalability experiments
  plot_results.py        turns results/*.csv into the figures used in this report
tests/
  test_eals_synthetic.py  correctness check on a small synthetic dataset
notebooks/
  eals_demo.ipynb           small-data, end-to-end, fully executed notebook
data/
  raw/                 downloaded Amazon ratings CSV
  processed/             filtered/reindexed train & test parquet
results/
  *.csv                    experiment outputs
  figures/*.png              figures used throughout this report
\end{lstlisting}

\section{Reproducing the experiments}

\begin{lstlisting}[language=bash, caption={Commands to reproduce every experiment in this report from scratch.}]
python -m venv .venv && source .venv/bin/activate
pip install pyspark numpy pandas scipy matplotlib jupyter nbformat pyarrow setuptools

# 1. Download and filter the dataset
curl -o data/raw/ratings_Movies_and_TV.csv \
  https://snap.stanford.edu/data/amazon/productGraph/categoryFiles/ratings_Movies_and_TV.csv
python src/preprocess.py --raw data/raw/ratings_Movies_and_TV.csv \
  --out data/processed --min-interactions 10

# 2. Run the weighting / convergence / scalability study
python src/run_experiments.py --data data/processed --out results
python src/plot_results.py --results results

# 3. Correctness check on synthetic data
python tests/test_eals_synthetic.py

# 4. Execute the small-data demonstration notebook end to end
jupyter nbconvert --to notebook --execute --inplace notebooks/eals_demo.ipynb
\end{lstlisting}

\section{The small-data demonstration notebook}
\label{app:notebook}

\texttt{notebooks/eals\_demo.ipynb} exercises the entire pipeline, from data
generation through training, evaluation, and inspection of the learned
model, on data small enough to read and verify by hand. It does not depend on the
large Amazon download being present (it is skipped gracefully, with a
printed message, if \texttt{data/processed/} does not exist). Each code
cell carries short comments explaining what it reads and what it produces
where that is not obvious from the code itself. The notebook was executed
top to bottom with \texttt{jupyter nbconvert --execute} immediately before
this report was finalized, with zero errors in any cell; its structure is:

\begin{enumerate}
  \item \textbf{Setup.} \emph{Input:} none. \emph{Output:} a local Spark
        session, with \texttt{src/eals.py} registered via
        \texttt{sc.addPyFile} so the closed-form update functions are
        importable on Spark's (local) executors, not just the driver.
  \item \textbf{Synthetic dataset.} \emph{Input:} none (a deterministic
        generator with a fixed seed). \emph{Output:} \texttt{train\_pairs}/
        \texttt{test\_pairs} for three disjoint 15-user/15-item
        communities (45 users, 45 items, 225 training interactions, 45
        held-out test interactions), small enough that a reader can
        manually verify the block structure the model is expected to
        recover.
  \item \textbf{Training.} \emph{Input:} the synthetic training pairs.
        \emph{Output:} a fitted \texttt{EALS} model
        (\texttt{model.P}, \texttt{model.Q}) and its per-iteration loss
        history, plotted to show monotonic-ish convergence (loss drops
        from 145.7 to 83.8 over 20 iterations).
  \item \textbf{Inspecting recommendations.} \emph{Input:} the trained
        model and a probed user id. \emph{Output:} that user's top-10
        recommended items, printed with their true community label. All
        ten come from the user's own community, the expected result
        for a correctly trained model on this synthetic structure.
  \item \textbf{Offline evaluation.} \emph{Input:} the trained model and
        the held-out test pairs. \emph{Output:} HR@10/NDCG@10 computed
        two ways: \texttt{evaluate\_full}, the paper's exact protocol
        (ranking against the whole 45-item catalog, tractable here because
        the catalog is tiny), and \texttt{evaluate\_sampled}, the
        relaxed protocol used for the large-scale experiments in this
        report.
  \item \textbf{Real-data subsample (optional).} \emph{Input:} a random
        300-user subsample of the preprocessed train/test parquet files
        under \texttt{data/processed/}, if present. \emph{Output:} the
        same training/evaluation pipeline
        run on real (if reduced) Amazon interaction data, producing
        HR@10/NDCG@10 for that subsample in well under a minute.
\end{enumerate}

\section{Full source listings}

The complete, unabridged source of every module described in
Chapter~\ref{ch:methods} is reproduced below, exactly as used to produce
the results in this report.

\subsection{\texttt{src/eals.py}: the eALS algorithm}
\lstinputlisting[language=Python]{../src/eals.py}

\subsection{\texttt{src/evaluate.py}: HR@K / NDCG@K evaluation}
\lstinputlisting[language=Python]{../src/evaluate.py}

\subsection{\texttt{src/baselines.py}: Spark MLlib ALS reference}
\lstinputlisting[language=Python]{../src/baselines.py}

\subsection{\texttt{src/preprocess.py}: dataset filtering and splitting}
\lstinputlisting[language=Python]{../src/preprocess.py}

\subsection{\texttt{src/run\_experiments.py}: experiment driver}
\lstinputlisting[language=Python]{../src/run_experiments.py}

\subsection{\texttt{src/plot\_results.py}: figure generation}
\lstinputlisting[language=Python]{../src/plot_results.py}

\subsection{\texttt{tests/test\_eals\_synthetic.py}: correctness check}
\lstinputlisting[language=Python]{../tests/test_eals_synthetic.py}
